\section{Methodology}
\label{sec:method}

\subsection{Audit Protocol}
\label{sec:audit_method}

For each dataset split, we record image count, storage format, resolution statistics, and file size. We treat two properties as likely shortcut sources:
\begin{itemize}[itemsep=2pt,topsep=2pt]
  \item \textbf{Format bias:} one class is predominantly JPEG and the other PNG.
  \item \textbf{Resolution bias:} one class occupies a narrow set of canonical sizes while the other spans many natural image sizes.
\end{itemize}

\subsection{Spectral Characterization}
\label{sec:spectral_method}

To quantify how much these construction choices separate the classes, we compute mean radial log-power spectra for real and fake images. Each image is converted to grayscale, transformed with a 2D FFT, and reduced to a 1D radial profile by azimuthal averaging~\cite{durall2020watch}. For each dataset we report a high-frequency divergence score,
\[
\mathrm{HF\text{-}L2} = \left\lVert \bar{S}_{\mathrm{real}}(r) - \bar{S}_{\mathrm{fake}}(r) \right\rVert_2
\quad \text{for } r > 0.8\,r_{\max},
\]
where $r_{\max}$ is the Nyquist radius. This is not meant to be a detector. It is a diagnostic for how strongly dataset construction separates the two classes before any model sees them.

\subsection{Detector Evaluation}
\label{sec:eval_method}

We evaluate six public detectors and one simple spectral CNN using their released checkpoints or the project's saved weights. Our core diagnostic is not only \emph{whether} AUC is high or low, but \emph{how it fails}.

\paragraph{Why below-chance AUC matters.}
If a detector lands near $0.5$ AUC on a new dataset, that may reflect ordinary out-of-distribution failure. If it lands well below $0.5$, then it is usually applying a decision rule in the wrong direction. In this paper we treat repeated AUC $< 0.5$ on the clean benchmark as evidence of an inverted rule, not just missing generalization.

\subsection{Normalization as a Diagnostic}
\label{sec:normalize}

To probe whether performance depends on file-format and size mismatches, we also evaluate a normalized variant of the test data. Every image is:
\begin{enumerate}[itemsep=1pt,topsep=2pt]
  \item Convert to RGB (strip alpha channel and EXIF metadata).
  \item Resize to $256 \times 256$ pixels using bilinear interpolation.
  \item Save as PNG (lossless; no re-JPEG-compression).
\end{enumerate}
This transform is applied identically to both classes. It does not remove content-baked artifacts already present in the pixels, but it does remove trivial differences in container format and nominal resolution. We therefore use normalization as a diagnostic, not as a claim that all confounds are solved.

\subsection{Format-Swap as a Causal Probe}
\label{sec:swap_method}

Normalization removes format differences but does not reverse them. To establish a causal link, we conduct a \emph{format-swap} experiment: every real image is re-encoded as PNG (losslessly from JPEG) and every fake image is re-encoded as JPEG at quality 85 (compressed from PNG). This reverses the natural format association present in both GenImage and the ProGAN training data that all six detectors were trained on.

If detector performance on GenImage relies on the JPEG$=$real/PNG$=$fake pattern, performance should collapse or invert when that pattern is reversed. We also run the swap on CIFAKE, where both classes are originally JPEG. Introducing a format difference on an otherwise artifact-free benchmark tests whether detectors will exploit format when given the opportunity.

\subsection{Frequency Band Ablation}
\label{sec:ablation_method}

To localize which frequency ranges carry the shortcut signal, we ablate images in the frequency domain before evaluation. Each image is transformed with a 2D FFT; all components within the target radial band are set to zero; the image is reconstructed via inverse FFT and clipped to $[0, 255]$. We test three non-overlapping bands defined by radial distance $r$ from the DC component:
\begin{itemize}[itemsep=1pt,topsep=2pt]
  \item \textbf{Low:} $r < 0.2\,r_{\max}$ --- DC, coarse structure, and format-induced energy differences.
  \item \textbf{Mid:} $0.2 \leq r < 0.6\,r_{\max}$ --- mid-range texture and edges.
  \item \textbf{High:} $r \geq 0.6\,r_{\max}$ --- fine detail, compression noise.
\end{itemize}
Ablation is applied to normalized $256\times256$ PNG images. A large AUC drop when a band is removed indicates detector reliance on that band's content.
